\documentclass{article}
\usepackage{amsmath,amssymb,amsthm}
\usepackage{algorithm}
\usepackage{algpseudocode}
\usepackage{geometry}
\geometry{margin=1in}
\newtheorem{theorem}{Theorem}
\newtheorem{lemma}{Lemma}
\newcommand{\E}{\mathbb{E}}

\begin{document}

\section*{Algorithm Name}
\textbf{Gradient Descent–Ascent for Convex–Concave Minimax
(GDA).}  
The method seeks a saddle point of a smooth function  
\[
\min_{y\in\mathcal Y}\;\max_{x\in\mathcal X}\;L(x,y),
\]
where $L:\mathcal X\times\mathcal Y\to\mathbb R$ is convex in $y$,
concave in $x$, and has Lipschitz continuous gradients.

-----------------------------------------------------------------------

\section*{Mathematical Formulation}
Let $z\coloneqq (x,y)\in\mathcal Z\equiv\mathcal X\times\mathcal Y$.
Define the \emph{monotone operator}
\[
F(z)\;:=\;
\begin{pmatrix}
-\nabla_x L(x,y)\\[2mm]
\;\;\nabla_y L(x,y)
\end{pmatrix}.
\tag{1}
\]
A saddle point $z^{\star}=(x^{\star},y^{\star})$ satisfies $F(z^{\star})=0$.
The GDA update with a (single) stepsize $\eta>0$ reads
\[
\boxed{
\begin{aligned}
x_{t+1}&=x_t+\eta\,\nabla_x L(x_t,y_t),\\[1mm]
y_{t+1}&=y_t-\eta\,\nabla_y L(x_t,y_t),
\end{aligned}}
\qquad t=0,1,2,\dots
\tag{2}
\]
or, compactly,
\[
z_{t+1}=z_t-\eta\,F(z_t).                                          \tag{3}
\]

-----------------------------------------------------------------------

\section*{Pseudocode}
\begin{algorithm}[H]
\caption{Gradient Descent–Ascent (GDA)}
\begin{algorithmic}[1]
\Require Initial point $z_0=(x_0,y_0)$, stepsize $\eta>0$, 
maximum iterations $T$.
\For{$t=0$ \textbf{to} $T-1$}
    \State Compute gradients $g^x_t\gets\nabla_x L(x_t,y_t)$,
          $g^y_t\gets\nabla_y L(x_t,y_t)$.
    \State $x_{t+1}\gets x_t+\eta\,g^x_t$ \Comment{ascent in $x$}
    \State $y_{t+1}\gets y_t-\eta\,g^y_t$ \Comment{descent in $y$}
\EndFor
\State \Return $z_T=(x_T,y_T)$ (or the averaged iterate
      $\bar z_T=\tfrac1T\sum_{t=1}^T z_t$).
\end{algorithmic}
\end{algorithm}

-----------------------------------------------------------------------

\section*{Dry Run Example}
We illustrate the update on a single‑variable quadratic
\[
f(w)= (w-3)^2 ,
\]
which is $L$-smooth with $L=2$.
Although $f$ is not a saddle‑point problem, the same
gradient‑descent step appears in (2) when $y$ is absent.
Take $w_0=0$, stepsize $\eta=0.1$, and run three iterations.

\begin{center}
\begin{tabular}{c|c|c|c}
$t$ & $w_t$ & $g_t=\nabla f(w_t)$ & $f(w_t)$\\ \hline
0 & $0$ & $2(w_0-3)=-6$ & $9$\\
1 & $w_1=w_0-\eta g_0=0-0.1(-6)=0.6$ &
    $2(0.6-3)=-4.8$ & $(0.6-3)^2=5.76$\\
2 & $w_2=w_1-\eta g_1=0.6-0.1(-4.8)=1.08$ &
    $2(1.08-3)=-3.84$ & $(1.08-3)^2=3.6864$\\
3 & $w_3=w_2-\eta g_2=1.08-0.1(-3.84)=1.464$ &
    $2(1.464-3)=-3.072$ & $(1.464-3)^2=2.3665$\\
\end{tabular}
\end{center}
The objective value decreases as expected,
illustrating the role of the gradient step that also appears
in the GDA updates (2).

-----------------------------------------------------------------------

\section*{Assumptions}
\begin{enumerate}
\item[(A1)] \textbf{Convex–concave structure.}  
For every $y$, $L(\cdot ,y)$ is concave; for every $x$, $L(x,\cdot )$ is convex.

\item[(A2)] \textbf{Smoothness.}  
There exists $L>0$ such that for all $z,z'\in\mathcal Z$,
\[
\|F(z)-F(z')\|\le L\|z-z'\| .                                           \tag{4}
\]

\item[(A3)] \textbf{Existence of a saddle point.}  
There is $z^{\star}=(x^{\star},y^{\star})\in\mathcal Z$ with $F(z^{\star})=0$.

\item[(A4)] \textbf{Bounded domain.}  
$\mathcal Z$ is convex and $\|z-z^{\star}\|\le D$ for all $z\in\mathcal Z$.
\end{enumerate}

-----------------------------------------------------------------------

\section*{Main Convergence Theorem}
\begin{theorem}[Ergodic convergence of GDA]\label{thm:main}
Assume (A1)–(A4) and choose a constant stepsize
\[
0<\eta\le\frac{1}{L}.                                                \tag{5}
\]
Let the averaged iterate be
\[
\bar z_T\;:=\;\frac1T\sum_{t=0}^{T-1}z_t .
\]
Define the primal–dual gap
\[
\mathcal G(z)\;:=\;
\max_{x\in\mathcal X}L(x,y)-\min_{y\in\mathcal Y}L(x,y).                \tag{6}
\]
Then the following bound holds:
\[
\boxed{
\frac1T\sum_{t=0}^{T-1}\mathcal G(z_t)
\;\le\;
\frac{D^{2}}{2\eta T}\;,
\qquad\text{and consequently}\qquad
\mathcal G(\bar z_T)\le\frac{D^{2}}{2\eta T}.}
\tag{7}
\]
Thus GDA attains an $O(1/T)$ rate in the ergodic sense.
\end{theorem}

-----------------------------------------------------------------------

\section*{Theoretical Properties}
\begin{itemize}
\item \textbf{Stability.}  
Inequality (7) shows that as long as $\eta\le1/L$, the
potential $\|z_t-z^{\star}\|^{2}$ is non‑increasing, guaranteeing
bounded iterates.

\item \textbf{Hyperparameter sensitivity.}  
If $\eta$ is chosen larger than $1/L$, the factor $(2-\eta L)$ in the
key descent inequality becomes negative and the method may diverge.
Hence the Lipschitz constant $L$ (or an upper bound on it) must be
known or estimated.

\item \textbf{Comparison to baselines.}  
Standard gradient descent on a convex problem also yields an
$O(1/T)$ ergodic rate under the same smoothness assumption.
Thus GDA does not lose asymptotic speed despite performing
ascent in $x$ and descent in $y$ simultaneously.
\end{itemize}

-----------------------------------------------------------------------

\section*{Preliminaries (Definitions and Lemmas)}
\begin{lemma}[Co‑coercivity of $L$‑smooth monotone operators]\label{lem:coerc}
If $F$ satisfies (4), then for all $z,z'\in\mathcal Z$
\[
\langle F(z)-F(z'),\,z-z'\rangle\;\ge\;\frac{1}{L}\,
\|F(z)-F(z')\|^{2}.                                             \tag{8}
\]
\end{lemma}
\begin{proof}
The $L$‑Lipschitz condition (4) implies
$\|F(z)-F(z')\|\le L\|z-z'\|$.
Apply the Cauchy–Schwarz inequality:
\[
\|F(z)-F(z')\|^{2}\le L\|F(z)-F(z')\|\|z-z'\|
\;\Longrightarrow\;
\langle F(z)-F(z'),z-z'\rangle\ge\frac{1}{L}\|F(z)-F(z')\|^{2}.
\]
\end{proof}

\begin{lemma}[Gap lower bound]\label{lem:gap}
For any $z=(x,y)$ and any saddle point $z^{\star}=(x^{\star},y^{\star})$,
\[
\langle F(z),\,z-z^{\star}\rangle\;\ge\;\mathcal G(z).               \tag{9}
\]
\end{lemma}
\begin{proof}
By definition of $F$,
\[
\langle F(z),z-z^{\star}\rangle
= -\langle\nabla_x L(x,y),x-x^{\star}\rangle
  +\langle\nabla_y L(x,y),y-y^{\star}\rangle .
\]
Using concavity of $L(\cdot ,y)$ and convexity of $L(x,\cdot )$,
\[
-\langle\nabla_x L(x,y),x-x^{\star}\rangle
   \ge L(x^{\star},y)-L(x,y),\qquad
\langle\nabla_y L(x,y),y-y^{\star}\rangle
   \ge L(x,y)-L(x,y^{\star}),
\]
and adding the two inequalities yields (9). \hfill\qed
\end{proof}

-----------------------------------------------------------------------

\section*{Main Proof (Step‑by‑Step Derivation)}
We now prove Theorem~\ref{thm:main}.  All derivations are
explicit; each non‑trivial manipulation is numbered.

\subsection*{Descent of the potential}
\begin{align}
\|z_{t+1}-z^{\star}\|^{2}
&= \bigl\|z_t-\eta F(z_t)-z^{\star}\bigr\|^{2}                                     \tag{Step 1}\\
&= \|z_t-z^{\star}\|^{2}
   -2\eta\langle F(z_t),z_t-z^{\star}\rangle
   +\eta^{2}\|F(z_t)\|^{2}.                                                       \tag{Step 2}
\end{align}
Because $F(z^{\star})=0$, monotonicity of $F$ implies
$\langle F(z_t)-F(z^{\star}),z_t-z^{\star}\rangle\ge0$, i.e.
\[
\langle F(z_t),z_t-z^{\star}\rangle\ge0.                                          \tag{Step 3}
\]

\subsection*{Bounding the gradient norm}
Apply Lemma~\ref{lem:coerc} with $z'=z^{\star}$:
\[
\langle F(z_t)-F(z^{\star}),z_t-z^{\star}\rangle
\ge\frac{1}{L}\|F(z_t)-F(z^{\star})\|^{2}
= \frac{1}{L}\|F(z_t)\|^{2}.                                                   \tag{Step 4}
\]
Rearranging gives
\[
\|F(z_t)\|^{2}\le L\,\langle F(z_t),z_t-z^{\star}\rangle .                     \tag{Step 5}
\]

\subsection*{One‑step contraction}
Insert (Step 5) into (Step 2):
\[
\|z_{t+1}-z^{\star}\|^{2}
\le \|z_t-z^{\star}\|^{2}
   -2\eta\langle F(z_t),z_t-z^{\star}\rangle
   +\eta^{2}L\,\langle F(z_t),z_t-z^{\star}\rangle .
\]
Factor the inner product:
\[
\|z_{t+1}-z^{\star}\|^{2}
\le \|z_t-z^{\star}\|^{2}
   -\eta\bigl(2-\eta L\bigr)\langle F(z_t),z_t-z^{\star}\rangle .               \tag{Step 6}
\]

\subsection*{Choice of stepsize}
Assume $\eta\le 1/L$.  Then $2-\eta L\ge1$, and (Step 6) simplifies to
\[
\|z_{t+1}-z^{\star}\|^{2}
\le \|z_t-z^{\star}\|^{2}
   -\eta\,\langle F(z_t),z_t-z^{\star}\rangle .                                 \tag{Step 7}
\]

\subsection*{Summation over iterations}
Sum (Step 7) for $t=0,\dots,T-1$:
\[
\|z_T-z^{\star}\|^{2}
\le \|z_0-z^{\star}\|^{2}
   -\eta\sum_{t=0}^{T-1}\langle F(z_t),z_t-z^{\star}\rangle .                 \tag{Step 8}
\]
Discard the non‑negative term $\|z_T-z^{\star}\|^{2}$ and use the bound
$\|z_0-z^{\star}\|\le D$ from (A4):
\[
\eta\sum_{t=0}^{T-1}\langle F(z_t),z_t-z^{\star}\rangle
\le D^{2}.                                                                     \tag{Step 9}
\]

\subsection*{From inner product to gap}
Lemma~\ref{lem:gap} (Step 9) yields $\langle F(z_t),z_t-z^{\star}\rangle
\ge\mathcal G(z_t)$.  Substituting into (Step 9) gives
\[
\sum_{t=0}^{T-1}\mathcal G(z_t)\;\le\;\frac{D^{2}}{\eta}.                     \tag{Step 10}
\]

\subsection*{Ergodic bound}
Dividing (Step 10) by $T$ we obtain the first inequality of (7):
\[
\frac1T\sum_{t=0}^{T-1}\mathcal G(z_t)\;\le\;\frac{D^{2}}{\eta T}.            \tag{Step 11}
\]
Finally, convexity of the gap functional implies
$\mathcal G(\bar z_T)\le\frac1T\sum_{t=0}^{T-1}\mathcal G(z_t)$,
which yields the second inequality of (7).  ∎

-----------------------------------------------------------------------

\section*{Rate Derivation (With Constants)}
From (Step 11) and the admissible stepsize (5) we have the explicit
rate
\[
\mathcal G(\bar z_T)\;\le\;\frac{D^{2}L}{2\,T},
\qquad\text{since we may set }\eta=\frac{1}{L}.
\tag{12}
\]
Thus the constant in front of $1/T$ is $D^{2}L/2$.

-----------------------------------------------------------------------

\section*{Special Cases}
\begin{enumerate}
\item \textbf{Strongly convex–strongly concave $L$.}  
If $L$ is $\mu$‑strongly convex in $y$ and $\mu$‑strongly concave in $x$,
the operator $F$ becomes $\mu$‑strongly monotone.
Repeating the proof with the stronger inequality
$\langle F(z_t),z_t-z^{\star}\rangle\ge\mu\|z_t-z^{\star}\|^{2}$
yields a linear convergence rate
\[
\|z_t-z^{\star}\|\le (1-\eta\mu)^{t}\|z_0-z^{\star}\|.
\]

\item \textbf{Stochastic gradients.}  
If only unbiased stochastic estimates $\tilde F(z_t)$ are available
with bounded variance $\E\|\tilde F(z_t)-F(z_t)\|^{2}\le\sigma^{2}$,
the same steps lead to an $O(1/\sqrt{T})$ rate for the expected gap,
exactly as in classical SGD (cf. Example 1).

\item \textbf{Different stepsizes for $x$ and $y$.}  
Let $\eta_x\neq\eta_y$.  The compact form becomes
$z_{t+1}=z_t-\tilde\eta\,\tilde F(z_t)$ with a diagonal scaling matrix.
The analysis proceeds with the induced norm
$\|z\|_{\tilde\eta}^{2}=z^{\top}\tilde\eta^{-1}z$ and yields essentially the
same $O(1/T)$ bound provided the two stepsizes satisfy
$\max\{\eta_x,\eta_y\}\le 1/L$.

\end{enumerate}

-----------------------------------------------------------------------

\section*{Discussion}
The proof above shows that the plain Gradient Descent–Ascent method,
under the standard smoothness and convex–concave assumptions,
converges with the optimal $O(1/T)$ ergodic rate without any
variance‑reduction or extrapolation tricks.  The analysis hinges only on
the monotonicity of the operator $F$ and the co‑coercivity implied by
smoothness.  Consequently, the result is robust to modest
modifications (e.g., different stepsizes, stochasticity) and provides a
baseline against which more sophisticated schemes such as Extra‑Gradient
or Optimistic Gradient can be compared.

\end{document}